% 洛伦兹群（综述）
% license CCBYSA3
% type Wiki

本文根据 CC-BY-SA 协议转载翻译自维基百科\href{https://en.wikipedia.org/wiki/Lorentz_group}{相关文章}

在物理学和数学中，洛伦兹群是指闵可夫斯基时空上的所有洛伦兹变换所构成的群。闵可夫斯基时空是所有（非引力）物理现象的经典与量子背景。洛伦兹群的名字来源于荷兰物理学家亨德里克·洛伦兹。

例如，以下定律、方程和理论都遵循洛伦兹对称性：
\begin{itemize}
\item 狭义相对论的运动学定律
\item 电磁学理论中的麦克斯韦场方程
\item 电子理论中的狄拉克方程
\item 粒子物理学的标准模型
\end{itemize}
洛伦兹群体现了自然界所有已知基本规律中时空的基本对称性。在时空区域足够小、引力差异可以忽略的情况下，物理规律与狭义相对论一样，保持洛伦兹不变性。
\subsection{基本性质}
洛伦兹群是庞加莱群（即闵可夫斯基时空所有等距变换所成的群）的一个子群。洛伦兹变换严格来说就是那些保持原点不变的等距变换。因此，洛伦兹群是闵可夫斯基时空等距群关于原点的各向同性子群。出于这个原因，洛伦兹群有时被称为齐次洛伦兹群，而庞加莱群则被称为非齐次洛伦兹群。洛伦兹变换是线性变换的一个例子；而闵可夫斯基时空的一般等距变换则是仿射变换。
\subsubsection{物理学定义}
\begin{figure}[ht]
\centering
\includegraphics[width=6cm]{./figures/c8d228c33aac897c.png}
\caption{} \label{fig_LuoLz_1}
\end{figure}
假设有两个惯性参考系 $(t, x, y, z)$ 与 $(t', x', y', z')$，以及两个点 $P_1, P_2$。洛伦兹群是所有在这两个参考系之间的变换集合，这些变换保持光在两点之间传播的速度不变：
\[
c^{2}(\Delta t')^{2}-(\Delta x')^{2}-(\Delta y')^{2}-(\Delta z')^{2}
=
c^{2}(\Delta t)^{2}-(\Delta x)^{2}-(\Delta y)^{2}-(\Delta z)^{2}.~
\]
用矩阵形式表示，这些变换是所有满足下式的线性变换 $\Lambda$：
\[
\Lambda^{\mathsf{T}} \eta \Lambda = \eta,
\qquad
\eta = \operatorname{diag}(1,-1,-1,-1).~
\]
\subsubsection{数学定义}
在数学上，洛伦兹群可以描述为不定正交群$O(1,3)$，即保持如下二次型的不变的矩阵李群：
\[
(t,x,y,z)\;\mapsto\; t^{2}-x^{2}-y^{2}-z^{2}~
\]
作用在$\mathbb{R}^4$上（赋予这一二次型的向量空间有时记作$\mathbb{R}^{1,3}$）。当写成矩阵形式时（参见经典正交群），这一二次型在物理学中被解释为闵可夫斯基时空的度规张量。
\subsubsection{符号说明}
在表示洛伦兹群时，$O(1,3)$与$O(3,1)$两种记号都被广泛使用。前者指保持一个符号型为$(+,-,-,-)$ 的度规的矩阵，后者则对应符号型为 $(-,+,+,+)$ 的度规。由于度规整体符号在定义方程中无关紧要，因此这两类矩阵群是相同的。

在当代，有部分领域倾向采用$(1,3)$的记号，而$(3,1)$的记号在现今实践和历史文献中依然大量存在。本文所描述的一切内容，对$O(3,1)$的记号同样成立，只需作相应的符号替换。类似的考虑同样适用于相关定义（例如$SO^{+}(1,3)$与$SO^{+}(3,1)$）。
